\documentclass[10pt,a4paper]{amsart}
\usepackage{amsmath,amssymb,amsthm,graphicx}
% For commutative diagrams we use Xy-pic
\usepackage[all]{xy}

% We use 2cell for 2-commutative diagrams.
\xyoption{2cell}
% Theorem environments.
%
\theoremstyle{plain}
\newtheorem{theorem}[subsection]{Theorem}
\newtheorem{proposition}[subsection]{Proposition}
\newtheorem{lemma}[subsection]{Lemma}

\theoremstyle{definition}
\newtheorem{definition}[subsection]{Definition}
\newtheorem{example}[subsection]{Example}
\newtheorem{exercise}[subsection]{Exercise}
\newtheorem{situation}[subsection]{Situation}

\theoremstyle{remark}
\newtheorem{remark}[subsection]{Remark}
\newtheorem{remarks}[subsection]{Remarks}

\numberwithin{equation}{subsection}

% Macros
%
\def\lim{\mathop{\mathrm{lim}}\nolimits}
\def\colim{\mathop{\mathrm{colim}}\nolimits}
\def\Spec{\mathop{\mathrm{Spec}}}
\def\Hom{\mathop{\mathrm{Hom}}\nolimits}
\def\Ext{\mathop{\mathrm{Ext}}\nolimits}
\def\SheafHom{\mathop{\mathcal{H}\!\mathit{om}}\nolimits}
\def\SheafExt{\mathop{\mathcal{E}\!\mathit{xt}}\nolimits}
\def\Sch{\mathit{Sch}}
\def\Mor{\mathop{\mathrm{Mor}}\nolimits}
\def\Ob{\mathop{\mathrm{Ob}}\nolimits}
\def\Sh{\mathop{\mathit{Sh}}\nolimits}
\def\NL{\mathop{N\!L}\nolimits}
\def\CH{\mathop{\mathrm{CH}}\nolimits}
\def\proetale{{pro\text{-}\acute{e}tale}}
\def\etale{{\acute{e}tale}}
\def\QCoh{\mathit{QCoh}}
\def\Ker{\mathop{\mathrm{Ker}}}
\def\Im{\mathop{\mathrm{Im}}}
\def\Coker{\mathop{\mathrm{Coker}}}
\def\Coim{\mathop{\mathrm{Coim}}}

% Boxtimes
%
\DeclareMathSymbol{\boxtimes}{\mathbin}{AMSa}{"02}

%
% Macros for moduli stacks/spaces
%
\def\QCohstack{\mathcal{QC}\!\mathit{oh}}
\def\Cohstack{\mathcal{C}\!\mathit{oh}}
\def\Spacesstack{\mathcal{S}\!\mathit{paces}}
\def\Quotfunctor{\mathrm{Quot}}
\def\Hilbfunctor{\mathrm{Hilb}}
\def\Curvesstack{\mathcal{C}\!\mathit{urves}}
\def\Polarizedstack{\mathcal{P}\!\mathit{olarized}}
\def\Complexesstack{\mathcal{C}\!\mathit{omplexes}}
% \Pic is the operator that assigns to X its picard group, usage \Pic(X)
% \Picardstack_{X/B} denotes the Picard stack of X over B
% \Picardfunctor_{X/B} denotes the Picard functor of X over B
\def\Pic{\mathop{\mathrm{Pic}}\nolimits}
\def\Picardstack{\mathcal{P}\!\mathit{ic}}
\def\Picardfunctor{\mathrm{Pic}}
\def\Deformationcategory{\mathcal{D}\!\mathit{ef}}
\usepackage{cleveref}
\crefname{originallemma}{original lemma}{original lemmas}
\Crefname{originallemma}{Original lemma}{Original lemmas}
\makeatletter
\newlabel{sites-cohomology-lemma-injective-abelian-sheaf-injective-presheaf}{{03F6}{source}{}{sourcepointer-1}{}}
\newlabel{sites-cohomology-lemma-injective-abelian-sheaf-injective-presheaf@cref}{{[originallemma][0][]03F6}{[1][0][]source}{}{}{}}
\makeatother
\crefname{originallemma}{original lemma}{original lemmas}
\Crefname{originallemma}{Original lemma}{Original lemmas}
\makeatletter
\newlabel{sites-cohomology-lemma-cech-spectral-sequence}{{03AZ}{source}{}{sourcepointer-2}{}}
\newlabel{sites-cohomology-lemma-cech-spectral-sequence@cref}{{[originallemma][0][]03AZ}{[1][0][]source}{}{}{}}
\makeatother
\title{The Cech-to-cohomology spectral sequence of an abelian sheaf on a site}
\author{The Stacks Project Authors}
\date{stacks-project:etale-cohomology.tex:local-snapshot-20261011; GNU FDL 1.2 or later}
\begin{document}
\maketitle
\noindent\textit{Editorial note.} Counted proposition: Tag 03OW, etale-cohomology.tex lines 1770-1785 with its own proof at lines 1787-1830 (statement immediately followed by its proof; no non-adjacent pairing).\par\medskip
\noindent\textit{Editorial note.} The context blocks reproduce the author's definition of the Cech complex (Tag 03OL) and the two standing results quoted inside the proof (Tags 03OR, 03OV, the latter together with its short proof). They are supporting material, not additional counted problems.\par\medskip
\noindent\textit{Editorial note.} The tag pointer 03AZ is a standard reference for the spectral sequence of a double complex; it is not used to replace any step of the counted proof.\par\medskip
\noindent\textit{Editorial note.} The retained preamble is the project's native preamble (Xy-pic setup and the theorem/macro block); documentclass amsart is supplied mechanically because the native preamble selects the project class file, which is not redistributed here.\par\medskip
\noindent\textit{Allowed standard prerequisites (editorial).}\begin{itemize}
\item Cech complex of a family of morphisms with fixed target. Stacks tag 03OL (retained context, source lines 1476-1497). Fixes the complex C\textasciicircum{}bullet(U, F) and the groups H\textasciicircum{}p(U, F) that the spectral sequence computes.
\item Injectives of abelian sheaves are injective as abelian presheaves. Stacks tag 03OV (retained context, source lines 1750-1768). Makes the rows of the double complex C\textasciicircum{}bullet(U, I\textasciicircum{}bullet) exact in positive degrees; used by the author at lines 1804-1808.
\item Vanishing of Cech cohomology of an injective presheaf in positive degrees. Stacks tag 03OR (retained context, source lines 1701-1706). Supplies the row exactness cited at line 1808.
\item Spectral sequence of a double complex. Stacks tag 03AZ (Cohomology on Sites, Lemma 03AZ). Standard reference cited by the author at lines 1826-1829 to read the E\_2-page off the double complex.
\end{itemize}\medskip
\ifdefined\hypertarget\hypertarget{sourcepointer-1}{}\fi
\noindent\textit{Source reference.} Original lemma 03F6: sources/stacks-etale-cohomology/tags.txt, source lines 4481--4481. Detail pointer for the retained context statement Tag 03OV, whose own proof is reproduced in the context block; cited by the author at etale-cohomology.tex line 1767.\par\medskip
\ifdefined\hypertarget\hypertarget{sourcepointer-2}{}\fi
\noindent\textit{Source reference.} Original lemma 03AZ: sources/stacks-etale-cohomology/tags.txt, source lines 4328--4328. Standard double-complex spectral sequence reference cited by the author at lines 1827-1829; the identification of the E\_2-page and of the abutment is retained in full in the counted proof.\par\medskip
\begin{definition}
\label{definition-cech-complex}
Let $\mathcal{C}$ be a category. Let
$\mathcal{U} = \{U_i \to U\}_{i \in I}$ be a family of morphisms of
$\mathcal{C}$ with fixed target. Assume that all the fibre products
$U_{i_0} \times_U \ldots \times_U U_{i_p}$ exists in $\mathcal{C}$.
Let $\mathcal{F} \in \textit{PAb}(\mathcal{C})$ be an abelian
presheaf. We define the {\it {\v C}ech complex}
$\check{\mathcal{C}}^\bullet(\mathcal{U}, \mathcal{F})$ by
$$
\prod_{i_0 \in I} \mathcal{F}(U_{i_0}) \to
\prod_{i_0, i_1 \in I} \mathcal{F}(U_{i_0} \times_U U_{i_1}) \to
\prod_{i_0, i_1, i_2 \in I}
\mathcal{F}(U_{i_0} \times_U U_{i_1} \times_U U_{i_2}) \to \ldots
$$
where the first term is in degree 0 and the maps are the usual ones.
The {\it {\v C}ech cohomology groups} are defined by
$$
\check{H}^p(\mathcal{U}, \mathcal{F}) =
H^p(\check{\mathcal{C}}^\bullet(\mathcal{U}, \mathcal{F})).
$$
\end{definition}

\begin{lemma}
\label{lemma-hom-injective}
Notation and assumptions as in Definition \ref{definition-cech-complex}.
If $\mathcal{I}$ is an injective object of $\textit{PAb}(\mathcal{C})$,
then $\check H^p(\mathcal{U}, \mathcal{I}) = 0$ for all $p > 0$.
\end{lemma}

\section{The {\v C}ech-to-cohomology spectral sequence}
\label{section-cech-ss}

\noindent
This spectral sequence is fundamental in proving foundational results on
cohomology of sheaves.

\begin{lemma}
\label{lemma-forget-injectives}
The forgetful functor $\textit{Ab}(\mathcal{C})\to \textit{PAb}(\mathcal{C})$
transforms injectives into injectives.
\end{lemma}

\begin{proof}
This is formal using the fact that the forgetful functor has a left adjoint,
namely sheafification, which is an exact functor. For more details see
Cohomology on Sites,
Lemma \ref{sites-cohomology-lemma-injective-abelian-sheaf-injective-presheaf}.
\end{proof}

\begin{theorem}
\label{theorem-cech-ss}
Let $\mathcal{C}$ be a site. For any covering
$\mathcal{U} = \{U_i \to U\}_{i \in I}$ of $U \in \Ob(\mathcal{C})$
and any abelian sheaf $\mathcal{F}$ on $\mathcal{C}$
there is a spectral sequence
$$
E_2^{p, q}
=
\check H^p(\mathcal{U}, \underline{H}^q(\mathcal{F}))
\Rightarrow
H^{p+q}(U, \mathcal{F}),
$$
where $\underline{H}^q(\mathcal{F})$ is the abelian presheaf
$V \mapsto H^q(V, \mathcal{F})$.
\end{theorem}

\begin{proof}
Choose an injective resolution $\mathcal{F}\to \mathcal{I}^\bullet$ in
$\textit{Ab}(\mathcal{C})$, and consider the double complex
$\check{\mathcal{C}}^\bullet(\mathcal{U}, \mathcal{I}^\bullet)$
and the maps
$$
\xymatrix{
\Gamma(U, I^\bullet) \ar[r] &
\check{\mathcal{C}}^\bullet(\mathcal{U}, \mathcal{I}^\bullet) \\
& \check{\mathcal{C}}^\bullet(\mathcal{U}, \mathcal{F}) \ar[u]
}
$$
Here the horizontal map is the natural map
$\Gamma(U, I^\bullet) \to
\check{\mathcal{C}}^0(\mathcal{U}, \mathcal{I}^\bullet)$
to the left column, and the vertical map is induced by
$\mathcal{F}\to \mathcal{I}^0$ and lands in the bottom row.
By assumption, $\mathcal{I}^\bullet$ is a complex of injectives in
$\textit{Ab}(\mathcal{C})$, hence by
Lemma \ref{lemma-forget-injectives}, it is a complex of injectives in
$\textit{PAb}(\mathcal{C})$. Thus, the rows of the double complex are
exact in positive degrees (Lemma \ref{lemma-hom-injective}), and
the kernel of $\check{\mathcal{C}}^0(\mathcal{U}, \mathcal{I}^\bullet)
\to \check{\mathcal{C}}^1(\mathcal{U}, \mathcal{I}^\bullet)$
is equal to
$\Gamma(U, \mathcal{I}^\bullet)$, since $\mathcal{I}^\bullet$
is a complex of sheaves. In particular, the cohomology of the total complex
is the standard
cohomology of the global sections functor $H^0(U, \mathcal{F})$.

\medskip\noindent
For the vertical direction, the $q$th cohomology group of the $p$th column is
$$
\prod_{i_0, \ldots, i_p}
H^q(U_{i_0} \times_U \ldots \times_U U_{i_p}, \mathcal{F})
=
\prod_{i_0, \ldots, i_p}
\underline{H}^q(\mathcal{F})(U_{i_0} \times_U \ldots \times_U U_{i_p})
$$
in the entry $E_1^{p, q}$. So this is a standard double complex spectral
sequence, and the $E_2$-page is as prescribed. For more details see
Cohomology on Sites,
Lemma \ref{sites-cohomology-lemma-cech-spectral-sequence}.
\end{proof}

\end{document}
